% Source: https://artofproblemsolving.com/community/c3h2511560_6_equation
% Stable identity: aops-topic-2511560-q01
\textbf{One variable is $0$.} If, say, $z=0$, the equation becomes $x+y=0$. So every triple with one entry $0$ and the other two opposite, $(t,-t,0)$ in any order with $t$ an integer, is a solution (this includes $(0,0,0)$).

\textbf{No variable is $0$.} Replacing $(x,y,z)$ by $(-x,-y,-z)$ changes the sign of both sides, so it maps solutions to solutions. Hence we may assume that at least two of the variables are positive.
\begin{itemize}
\item \emph{All three positive.} By symmetry let $x\le y\le z$. Then $xyz=x+y+z\le3z$, so $xy\le3$, and $(x,y)$ is $(1,1)$, $(1,2)$ or $(1,3)$. The first gives $2+z=z$, impossible; the second gives $3+z=2z$, so $z=3$; the third gives $4+z=3z$, so $z=2<y$, contradicting $y\le z$. So $\{x,y,z\}=\{1,2,3\}$, and indeed $1+2+3=6=1\cdot2\cdot3$.
\item \emph{Exactly two positive}, say $x,y>0>z$. Put $z=-w$ with $w>0$: then $x+y-w=-xyw$, i.e. $x+y=w(1-xy)$. The left side is positive, while $xy\ge1$ makes the right side at most $0$. No solutions.
\end{itemize}
Undoing the sign change adds the triples formed by $-1,-2,-3$.

\textbf{Answer.} The solutions are all permutations of $(1,2,3)$ and of $(-1,-2,-3)$, together with all permutations of $(t,-t,0)$ for integers $t$:
\[\boxed{(x,y,z)\in\{\text{perm. of }(1,2,3)\}\cup\{\text{perm. of }(-1,-2,-3)\}\cup\{\text{perm. of }(t,-t,0):t\in\mathbb{Z}\}}\]
